\documentclass{article}
\usepackage[margin=2cm]{geometry}
\usepackage{amsmath}
\usepackage{graphicx}
\usepackage{booktabs}
\usepackage[utf8]{inputenc}
\usepackage{ragged2e}
\setlength{\emergencystretch}{3em}
\usepackage{paracol}
\begin{document}
\sloppy

\noindent\hfill\includegraphics[height=1.8em]{assets/figure_header_logo.png}\par\noindent\hrule\vspace{4pt}

\begin{paracol}{2}
\RaggedRight

However, this approach quickly becomes infeasible as it
requires an exponential number of models with respect to
the dimensionality of the input space. In this work, we propose several conditional transformations that handle arbitrary dimensionality in a principled manner. Our contributions are as follows. 1) We propose a novel extension of
flow-based generative models to model the conditional distribution of arbitrary unobserved covariates in data imputation tasks. Our method is the first to develop invertible
transformations that operate on an arbitrary set of covariates. 2) We strengthen a flow-based model by using a novel
autoregressive conditional likelihood. 3) We propose a novel
penalty to generate a single imputed "best guess for models without an analytically available mean. 4) We run extensive empirical studies and show that achieves state-of-the-art
performance for both missing feature imputation and image
inpainting on benchmark real-world datasets.


\subsection*{1Problem Formulation}


Consider a real-valued distribution $p ( { \pmb x } )$ over $\mathbb { R } ^ { \mathbf { d } }$ We are
interested in estimating the conditional distribution of all
possible subsets of covariates $u \subseteq \{ 1, \dots, d \}$ conditioned
on the remaining observed covariates${ \mathfrak { o } } ~ = ~ \{ 1, \dots, d \} ~ { \mathfrak { u } }.$
That is, we shall estimate $\pmb { p } ( \pmb { x } . { \pmb { u } } ~ \mid \pmb { x } . { o } )$ where $\pmb { x } . { \pmb { u } } ~ \in ~ \mathbb { R } ^ { | \pmb { u } | }$
and$\pmb { x } . { \pmb { a } } ~ \in ~ \mathbb { R } ^ { | \pmb { a } | }$ for all possible subsets.$^ { u }$ For ease of notation, let$b \in \{ 0, 1 \} ^ { a }$ be a binary mask indicating which dimensions are observed. Furthermore, let${ \pmb \nu } [ { \mathfrak { b } } ]$index dimensions$\imath . { \mathbf { i } }$ for which $\mathfrak { b } . { \mathrm { i } } ~ = ~ 1$ for the bitmask $^ { h }$ on a vector $t ^ { 3 }$ Thus, $x . { a } ~ = ~ x \left [ { \mathfrak { h } } \right ]$ denotes observed dimensions and
$\pmb { x } . { \pmb { u } } = \pmb { x } \left [ 1 - \pmb { b } \right ]$denotes unobserved dimensions. We also apply this index to matrices such that $M [ b, b ]$ indexes rows
and then columns. Without loss of generanty, conditionals
may also be conditioned on the bitmask $b$ $p ( \mathbf { r } . { u } \mid \mathbf { r } . { o }, \mathbf { h } )$
and will be estimated with maximum log-likelihood estimation as described below. In addition, imputation tasks shall
be accomplished by generating samples from the conditional
distributions${ \mathfrak { p } } ( { \mathfrak { x } } . { \mathfrak { u } } \mid { \mathfrak { x } } . { \mathfrak { o } }, { \mathfrak { b } } )$


\subsection*{2Background}


builds on Transformation Autoregressive Networks (TANs)
, a flow-based model that combines transformation of vari-.
ables with autoregressive likelihoods. We expound on flowbased models and TANs below. The change of variable theorem (??) is the cornerstone of flow-based generative models,
where represents an invertible transformation that trans-$q$
forms covariates from input space$\boldsymbol { \mathcal { X } } ^ { \prime }$into a latent space Z.


\begin{equation}
\begin{array} { r } { \underbrace { \cdots \cdots \cdots \cdots \cdots \cdots \cdots \cdots \cdots \cdots \cdots \cdots \cdots \cdots \cdots \cdots \cdots \cdots \cdots \cdots \cdots \cdots \cdots \cdots \cdots \cdots \cdots \cdots \cdots } _ { ( 1 ) } } \\ { p _ { \mathcal { X } } ( x ) = \left. \underbrace { p _ { \mathcal { X } } ( q ( x ) ) } _ { \mathrm { ~ t o n i s ~ e r a m e d ~ t h a r e r a n d ~ s t h a ~ t r e r a n d ~ f o r e r e } } \right. } \end{array}
\end{equation}


In order to efficiently compute the determinant, the transformation is often designed to have a diagonal or triangular$t . { t }$
tive, the flow models are often composed of multiple transformations in a sequence to get a more flexible composite
transformation, i.e. $q = q . { m } \circ q . { m - 1 } \circ \dots \circ q . { 1 }$ Here, the covariates flow through a chain of transformations, substituting
the last output variable as input for the next transformation.
Typically, a flow-based model transforms the covariates to a

\switchcolumn
\RaggedRight

latent space with a simple base distribution, like a standard
Gaussian. However, TANs provides more flexibility by mod
eling the latent distribution with an autoregressive approach
 This alters the earlier equation, in that ${ \mathfrak { p } } { \mathfrak { z } } \{ { \mathfrak { q } } ( { \mathfrak { x } } ) \}$ ISnow
represented as the product of  conditional distributions.$d$


\begin{equation}
\begin{array} { r } { { p _ { \mathrm { X a n c e ~ m o r ~ ~ s u r e ~ ~ c o n ~ } } } ~ }\end{array}
\end{equation}


Since flow models give the exact likelihood, they can be
trained by directly optimizing the log likelihood. In addition, thanks to the invertibility of the transformations, one
can draw samples by simply inverting the transformations
over a set of samples from the latent space.


\subsection*{3Methods}


We now develop by constructing both conditional trans-.
formations of variables and autoregressive likelihoods that
work with an arbitrary set of unobserved covariates. To deal
with arbitrary dimensionality for conditioning covariates$x . { \alpha }$
we define a zero imputing function ${ \dot { \phi } } ( r : h )$ that imputes input vectorwith zeros based on the specitied binary mask$\mathbf { v }$
b:


\begin{equation}
\begin{array} { c } { w = \phi ( v ; b ), \quad w _ { i } = \left \{ \begin{array} { l l } { v _ { c }, } & { \quad b _ { i } = 1 } \\ { 0, } & { \quad b _ { i } = 0 } \end{array} \right., \qquad ( 3 ) } \\ { \cdots ~, \quad ~ }}\end{array}
\end{equation}


where ${ \mathfrak { c } } ~ = ~ \sum . { j = 1 } ^ { \mathrm { i } } b . { j }$ represents the cumulative sum over
$^ { h }$ For example, ${ \dot { \varphi } } \{ { \pmb x } . { \alpha } ; { \pmb b } \}$gives adimensional vector with$d$
tion.). Thus, we get a conditioning vector with fixed dimensionality. However, handling the arbitrary dimensionality of
$x . { u }$ requires further care, as discussed below. We first consider a conditional extension to the change of variable theorem:


\begin{equation}
\begin{array} { r l } & { \mathrm { t e n a, ~ } } \\ & { p _ { \mathcal { X } } ( x _ { u } \mid x _ { o }, b ) = \left | \operatorname* { d e t } \frac { d q _ { x _ { u }, b } } { d x _ { u } } \right | p _ { \mathcal { Z } } ( q _ { x _ { u }, b } ( x _ { u } ) \mid x _ { o }, b ), ~ ~ ( 4 ) } \\ & { \mathrm { w h e r m ~ o ~ } \quad, ~ ~ \mathrm { i s e ~ a ~ t r m}}\end{array}
\end{equation}


where $q . { a, b }$ is a transformation on the unobserved covariates$x . { \mathbf { u } }$with respect to the observed covariates$\mathcal { X } . { a }$ and binary mask as demonstrated in Fig.. However, the fact that$b$
$x . { \mathbf { u } }$
ing dimensions makes it challenging to define $q . { x . { o }, b }$ s across
different bitmasks $b$
the transformation to have adaptive outputs that can adapt
to different dimensionality of$x . { u }$ Another challenge is that
 different missing patterns require the transformation to capture different dependencies. Since the missing pattern could
be arbitrary, we require the transformation to learn a large
range of possible dependencies. Aiming at solving those
challenges, we propose several conditional transformations
that leverage the conditioning information in$x . { u }$and $b$ and
can be adapted to$x . { \mathbf { u } }$with arbitrary dimensionality. We describe them in detail below. For notation simplicity, we drop
the subscripts$\boldsymbol { \chi } ^ { * }$ $z$in the following sections.


Affine Coupling TransformationAffine couplingis a
commonly used flow transformation. Here, we derive a conditional extension to it. Just as in the unconditional counterparts , we divide the unobserved covariates$r . { u }$ into two
parts

\end{paracol}

\end{document}
